% p028_a 的电脑重画（基尔霍夫电压定律例图：电池+开关供电；电流表 A、R1、R2(变阻器)；电压表 V1 并联 R1、V2 并联 R2、左侧 V2 接在电源两端；电压表旁箭头向下；大椭圆为“找圈”标记）
\begin{tikzpicture}[line width=0.7pt, >={Stealth[length=2.2mm]}]
  % ---- 电池与开关支路（左） ----
  \draw (0,5) -- (0,4.5);
  \draw (0,4.45) circle (0.05);
  \draw (0.03,4.41) -- (0.3,3.95);
  \draw (0,3.9) -- (0,2.4);
  \draw[<-, >={Stealth[length=1.6mm]}] (-0.5,2.4) -- (0.5,2.4);
  \draw (-0.25,2.1) -- (0.25,2.1);
  \draw (0,2.1) -- (0,0) -- (4.8,0);
  % ---- 上边：电流表 ----
  \draw (0,5) -- (2.67,5);
  \node[draw, circle, inner sep=0pt, minimum size=0.66cm] at (3,5) {A};
  \draw (3.33,5) -- (4.8,5) -- (4.8,4.08);
  % ---- 左侧电压表 V2 ----
  \draw (1.6,5) -- (1.6,2.78);
  \node[draw, circle, inner sep=0pt, minimum size=0.76cm] at (1.6,2.4) {V$_2$};
  \draw (1.6,2.02) -- (1.6,0);
  \draw[->] (2.2,3.1) -- (2.2,1.7);
  % ---- R1、R2 支路 ----
  \draw (3.6,5) -- (3.6,3.85);
  \draw (3.42,2.95) rectangle (3.78,3.85);
  \draw (3.6,2.95) -- (3.6,2.15);
  \draw (3.42,1.25) rectangle (3.78,2.15);
  \draw (3.6,1.25) -- (3.6,0);
  \draw (3.42,1.85) -- (3.1,1.85) -- (3.1,0.92) -- (3.6,0.92);
  \node at (3.0,3.4) {$R_1$};
  \node at (2.65,1.7) {$R_2$};
  % ---- 右侧电压表 V1（并联 R1）、V2（并联 R2） ----
  \draw (3.6,2.5) -- (4.8,2.5);
  \node[draw, circle, inner sep=0pt, minimum size=0.76cm] at (4.8,3.7) {V$_1$};
  \draw (4.8,3.32) -- (4.8,1.58);
  \node[draw, circle, inner sep=0pt, minimum size=0.76cm] at (4.8,1.2) {V$_2$};
  \draw (4.8,0.82) -- (4.8,0);
  \draw[->] (5.4,4.4) -- (5.4,3.0);
  \draw[->] (5.4,1.9) -- (5.4,0.5);
  % ---- 结点 ----
  \fill (1.6,5) circle (1.6pt) (3.6,5) circle (1.6pt) (3.6,2.5) circle (1.6pt)
        (4.8,2.5) circle (1.6pt) (1.6,0) circle (1.6pt) (3.6,0) circle (1.6pt);
  % ---- 大椭圆（找圈） ----
  \draw (3.4,2.3) ellipse (2.75 and 3.55);
\end{tikzpicture}
